\subsection{Why Can I Know the Right Move and Still Fail to Start?}

The short answer is that knowledge is not the same as transition power. You can
understand the target state perfectly and still remain stuck if framework
inertia is still strong, the current barrier is higher than your available
volitional energy, or the decision window has already narrowed before any
stabilizing action begins. In practice, the solution is not to argue with
yourself longer. It is to make the first move smaller, clearer, and earlier.

\begin{enumerate}[nosep]
  \item Name the target state in one sentence, without extra explanation.
  \item Identify the current state that is keeping the old frame alive.
  \item Find the smallest barrier term you can change now: ambiguity, setup
        cost, missing tools, or location mismatch.
  \item Open a decision window by choosing a concrete time and place for the
        first move.
  \item Reduce the first move to a visible action that takes less volitional
        energy than the full task.
  \item Stabilize the new state immediately after the first move so the old
        frame does not rebuild itself.
\end{enumerate}

\begin{remark}
Suppose you know you should write the report, but you still cannot begin.
The useful question is not ``Why am I lazy?'' The useful question is ``What is
the first action that lowers the barrier enough to let the new frame start?''
If the answer is ``open the file, write the heading, and sit at the desk,''
then that is the correct beginning. The value of this approach is that it
converts a vague intention into a designed transition instead of a private
moral test.
\end{remark}
